\chapter{System Model and Problem Formulation}
\label{ch:system-model}

This chapter defines the system model used throughout the thesis. The objective is
to provide a consistent notation for the vehicular network, the multimodal sensing
observations, the communication beam measurements, the target-user association
variables, and the simplified communication-performance abstraction. The model is
intended to support the sensing-assisted beam management framework developed in
the following chapter; it is therefore kept at a system level and avoids assuming
implementation details that have not yet been fixed.

\section{Network Scenario and Time Structure}
\label{sec:system-network-time}

We consider a roadside vehicular network consisting of base stations (BSs),
distributed sensing nodes (SNs), and mobile vehicle users. Let
\(\mathcal{B}\) denote the set of BSs. A BS provides downlink communication
service to vehicles in its local coverage region and also acts as the center of a
local sensing-assisted beam management process. Let \(\mathcal{S}_b\) denote the
set of SNs selected by BS \(b \in \mathcal{B}\) for cooperative sensing. The
number of cooperating SNs is denoted by \(K_{\mathrm{SN}}=|\mathcal{S}_b|\), and
is treated as a configurable system parameter rather than a fixed constant in
this first-pass formulation. Each SN contributes sensing observations but does
not directly provide communication service in the considered model.

The mobile communication users are vehicles traveling through the road region.
For a serving BS \(b\), the set of active vehicle users at decision epoch \(t\)
is denoted by \(\mathcal{U}_b(t)\). Each communication user corresponds to a
physical vehicle target in the sensing scene, but the sensing target identity and
the communication user identity are not assumed to be known a priori. This
distinction is central to the target-user association problem addressed in this
thesis.

Time is divided into discrete decision epochs \(t=0,1,\ldots,T-1\). Each epoch
may correspond to one sensing frame, one communication scheduling interval, or a
coarser beam-management interval depending on the final simulation setup. At
each epoch, BS \(b\) receives multimodal sensing observations from its local
sensor suite and from the cooperating SNs, maintains tracking states for
previously associated users, and selects a beam-management action for the active
vehicle users. Initial access epochs may include a more complete beam sweep,
whereas subsequent epochs are intended to rely more heavily on sensing-assisted
prediction and tracking.

Each BS and each SN is equipped with four tilted down-looking cameras that cover
the road region from multiple viewpoints. In addition, each BS has an ISAC
sensing capability that produces a local sensing point cloud. For BS \(b\), all
sensing inputs are represented in a local bird's-eye-view (BEV) coordinate frame
anchored to the BS or to a fixed road-side reference frame. The exact simulation
layout, the default value of \(K_{\mathrm{SN}}\), the frame duration, and the
spatial extent of the BEV region remain configuration choices.

TODO[unstable-method]: Fix the final BS/SN layout, the default number of
cooperative SNs, the frame duration, and the BEV coordinate range after the
simulation system is finalized.

\section{Multimodal Sensing Observation Model}
\label{sec:system-sensing-observation}

At epoch \(t\), BS \(b\) observes the road scene through camera and ISAC
modalities. The multiview camera observations at the BS are denoted by
\[
    \mathcal{I}_b(t)=\{I_{b,c}(t)\}_{c=1}^{4},
\]
where \(I_{b,c}(t)\) is the image captured by the \(c\)-th camera mounted on BS
\(b\). Similarly, the camera observations from the cooperating SNs are denoted
by
\[
    \mathcal{I}_{\mathcal{S}_b}(t)
    = \{I_{s,c}(t): s \in \mathcal{S}_b,\ c=1,\ldots,4\}.
\]
The BS-side ISAC modality is represented by a point-cloud observation
\[
    \mathcal{P}_b(t)=\{(\mathbf{q}_{b,m}(t), a_{b,m}(t))\}_{m=1}^{N_b^{\mathrm{pt}}(t)},
\]
where \(\mathbf{q}_{b,m}(t)\) denotes the three-dimensional location of the
\(m\)-th sensing point in the local coordinate system, \(a_{b,m}(t)\) denotes an
optional point attribute such as intensity or confidence, and
\(N_b^{\mathrm{pt}}(t)\) is the number of points available at epoch \(t\). This
point cloud is used as an
ISAC-like geometric observation and should not be interpreted as a fully
specified radar signal-processing model.

The complete sensing observation available to BS \(b\) is written as
\begin{equation}
    \mathcal{O}^{\mathrm{sen}}_b(t)
    = \left\{\mathcal{I}_b(t),\mathcal{I}_{\mathcal{S}_b}(t),\mathcal{P}_b(t)\right\}.
    \label{eq:sensing-observation}
\end{equation}
The camera images and point cloud are assumed to be temporally synchronized and
spatially calibrated before being fused. Let \(\mathcal{T}_{j \rightarrow b}\)
denote the calibrated transformation from the coordinate frame of sensor or node
\(j\) into the local BEV frame of BS \(b\). These transformations allow visual
features and ISAC point-cloud features to be encoded into a common BEV
representation. The detailed projection, lifting, and point-cloud encoding
operations are left to the method chapter, because their exact implementation is
still part of the model design.

After multimodal BEV fusion, a vehicle detection module outputs a set of detected
targets
\[
    \mathcal{D}_b(t)=\{d_{b,1}(t),\ldots,d_{b,N_b^{\mathrm{det}}(t)}(t)\}.
\]
For detected target \(d_{b,i}(t)\), the sensing-side target representation is
summarized as
\begin{equation}
    \mathbf{z}_{b,i}(t)
    = \left(\hat{\mathbf{p}}_{b,i}(t),\hat{\mathbf{l}}_{b,i}(t),
    \hat{c}_{b,i}(t)\right),
    \label{eq:target-representation}
\end{equation}
where \(\hat{\mathbf{p}}_{b,i}(t)=[\hat{x}_{b,i}(t),\hat{y}_{b,i}(t)]^{\mathrm{T}}\)
is the estimated BEV center, \(\hat{\mathbf{l}}_{b,i}(t)\) contains target-size
or box-related attributes when available, and \(\hat{c}_{b,i}(t)\) is a detection
confidence score. Additional attributes such as heading or appearance features
may be included if they are used by the final association and tracking modules.

TODO[unstable-method]: Specify the final BEV resolution, feature encoding,
camera-to-BEV projection, ISAC point-cloud representation, and detection-target
attributes after the network architecture is fixed.

\section{Communication and Beam Model}
\label{sec:system-communication-beam}

Each BS is equipped with a finite beam codebook for directional communication.
For BS \(b\), the codebook is denoted by
\[
    \mathcal{F}_b=\{\mathbf{f}_{b,1},\mathbf{f}_{b,2},\ldots,\mathbf{f}_{b,M_b^{\mathrm{beam}}}\},
\]
where \(\mathbf{f}_{b,m}\) is the \(m\)-th beamforming vector and
\(M_b^{\mathrm{beam}}\) is the codebook size. For a vehicle user
\(u \in \mathcal{U}_b(t)\), the channel from BS \(b\) at epoch \(t\) is denoted
by \(\mathbf{h}_{b,u}(t)\). The idealized beam quality of codebook beam \(m\) can
be expressed as
\begin{equation}
    g_{b,u,m}(t)
    = \left|\mathbf{h}_{b,u}^{\mathrm{H}}(t)\mathbf{f}_{b,m}\right|^2,
    \label{eq:beam-quality}
\end{equation}
up to a transmit-power and noise normalization. The vector
\(\mathbf{g}_{b,u}(t)=[g_{b,u,1}(t),\ldots,g_{b,u,M_b^{\mathrm{beam}}}(t)]\)
is referred to as the beam power distribution of user \(u\) with respect to the
codebook of BS \(b\).

The sensing-assisted beam prediction module does not directly know
\(\mathbf{g}_{b,u}(t)\), because its input is a detected sensing target rather
than a communication identity. For each detected target \(d_{b,i}(t)\), it
predicts a per-target beam quality vector
\begin{equation}
    \hat{\mathbf{g}}_{b,i}(t)
    = [\hat{g}_{b,i,1}(t),\ldots,\hat{g}_{b,i,M_b^{\mathrm{beam}}}(t)].
    \label{eq:predicted-beam-distribution}
\end{equation}
The predicted Top-\(K\) candidate beam set is then
\begin{equation}
    \hat{\mathcal{K}}_{b,i}(t)
    = \mathrm{TopK}_{K_{\mathrm{beam}}}\left(\hat{\mathbf{g}}_{b,i}(t)\right),
    \label{eq:topk-beams}
\end{equation}
where \(K_{\mathrm{beam}}\) denotes the number of candidate communication beams
retained for later beam management and evaluation. The notation
\(\mathrm{TopK}_{K_{\mathrm{beam}}}(\cdot)\) returns the indices of the
\(K_{\mathrm{beam}}\) largest entries in its argument.

During initial access, BS \(b\) may perform an SSB-like beam sweep for a
communication user \(u\). The resulting measurement is abstracted as a beam power
spectrum
\begin{equation}
    \mathbf{y}^{\mathrm{SSB}}_{b,u}(t_0)
    = [y^{\mathrm{SSB}}_{b,u,1}(t_0),\ldots,
       y^{\mathrm{SSB}}_{b,u,M_b^{\mathrm{beam}}}(t_0)],
    \label{eq:ssb-spectrum}
\end{equation}
where \(t_0\) is an initial access or re-association epoch. This measurement is
used to establish the correspondence between communication users and detected
sensing targets. In later epochs, the system may use only partial beam
measurements or calibration signals over a smaller candidate set, for example
over the predicted Top-\(K\) beams. Such measurements are modeled as auxiliary
communication observations rather than as a complete 5G NR protocol simulation.

This codebook-level abstraction is consistent with sensing-aided beam prediction
studies that supervise beam indices, Top-\(K\) candidate beams, or per-beam
quality distributions rather than modeling the complete standard-compliant
control procedure \cite{charan2021visionpositionmultimodalbeamprediction,alkhateeb2023deepsense6glargescalerealworld,charan2023camerabasedmmwavebeam,zeng2026bevfusionbasedframeworksequential}.

TODO[unstable-method]: Fix the final codebook size, beam-label generation
procedure, beam-power normalization, and abstraction of partial measurement or
calibration signals.

\section{Target-User Association and Tracking State}
\label{sec:system-association-tracking}

The system must connect three types of identity: detected sensing targets,
communication users, and temporal tracks. A detected target \(d_{b,i}(t)\)
describes an object observed in the BEV sensing output at epoch \(t\). A
communication user \(u\) is the network identity used for beam measurement,
scheduling, and throughput evaluation. A track \(r\) is a temporal state that
maintains the estimated motion of a physical vehicle across epochs.

Let \(a_{b,i,u}(t) \in \{0,1\}\) be the binary target-user association variable:
\begin{equation}
    a_{b,i,u}(t)=
    \begin{cases}
    1, & \text{if detected target } d_{b,i}(t) \text{ is associated with user } u,\\
    0, & \text{otherwise.}
    \end{cases}
    \label{eq:association-variable}
\end{equation}
The association is subject to one-to-one matching constraints at a given epoch,
\begin{equation}
    \sum_{u \in \mathcal{U}_b(t)} a_{b,i,u}(t) \leq 1,\quad
    \sum_{i=1}^{N_b^{\mathrm{det}}(t)} a_{b,i,u}(t) \leq 1.
    \label{eq:association-constraints}
\end{equation}
At an initial access epoch \(t_0\), the association may be estimated by matching
the sensing-predicted beam distribution \(\hat{\mathbf{g}}_{b,i}(t_0)\) with the
communication-side beam spectrum \(\mathbf{y}^{\mathrm{SSB}}_{b,u}(t_0)\), while
also considering BEV spatial priors and detection confidence when available. A
generic initial association problem can be written as
\begin{equation}
    \min_{\{a_{b,i,u}(t_0)\}}
    \sum_{i=1}^{N_b^{\mathrm{det}}(t_0)}
    \sum_{u \in \mathcal{U}_b(t_0)}
    a_{b,i,u}(t_0) C^{\mathrm{init}}_{b,i,u}(t_0),
    \label{eq:initial-association}
\end{equation}
subject to \eqref{eq:association-constraints}, where
\(C^{\mathrm{init}}_{b,i,u}(t_0)\) is an initial matching cost. In this thesis,
the cost is expected to include beam-spectrum similarity and may include
additional BEV-domain terms, but its final form is method-dependent.

For each associated vehicle, the system maintains a tracking state in the BEV
plane. Let \(r\) denote a track identity associated with a physical vehicle. A
minimal kinematic state is
\begin{equation}
    \mathbf{x}_{b,r}(t)
    = [x_{b,r}(t),y_{b,r}(t),v_{x,b,r}(t),v_{y,b,r}(t)]^{\mathrm{T}},
    \label{eq:tracking-state}
\end{equation}
where \((x_{b,r}(t),y_{b,r}(t))\) is the BEV position and
\((v_{x,b,r}(t),v_{y,b,r}(t))\) is the planar velocity. To support maneuvering
vehicles, the tracker may use an interacting multiple model (IMM) filter with a
set of motion models \(\mathcal{Q}\). The probability of motion model
\(q \in \mathcal{Q}\) for track \(r\) is denoted by \(\mu_{b,r,q}(t)\), with
\begin{equation}
    \sum_{q \in \mathcal{Q}} \mu_{b,r,q}(t)=1.
    \label{eq:imm-mode-probability}
\end{equation}
The IMM prediction provides a prior track state at epoch \(t\), which is then
associated with the new BEV detections \(\mathcal{D}_b(t)\). A generic
tracking-to-detection cost may be written as
\begin{equation}
    C^{\mathrm{trk}}_{b,r,i}(t)
    = \Psi_{\mathrm{trk}}\left(\hat{\mathbf{x}}^{-}_{b,r}(t),
    \mathbf{z}_{b,i}(t),\hat{\mathbf{g}}_{b,i}(t)\right),
    \label{eq:tracking-cost}
\end{equation}
where \(\hat{\mathbf{x}}^{-}_{b,r}(t)\) is the predicted track state before
measurement update and \(\Psi_{\mathrm{trk}}(\cdot)\) denotes a method-specific
association function. After a successful track-to-detection association, the
track state, the target-user mapping, and the beam prediction for the
corresponding communication user are updated jointly.

TODO[unstable-method]: Specify the final association cost, matching algorithm,
IMM motion models, gating thresholds, failure recovery rule, and whether beam
prediction is used directly in temporal association.

\section{Communication Performance Abstraction}
\label{sec:system-performance-abstraction}

The communication evaluation in this thesis uses a simplified and interpretable
system-level abstraction. It connects sensing-assisted Top-\(K\) beam prediction,
target-user association, tracking errors, and beam-management overhead to
multi-user communication metrics. It is not intended to reproduce the full 5G NR
protocol stack, detailed control-channel design, or exact standard-compliant
SSB/CSI-RS timing.

Let \(m_{b,u}(t)\) denote the selected beam index for user \(u\) served by BS
\(b\) at epoch \(t\). If the true effective beam for user \(u\) is contained in
the predicted candidate set associated with that user, the system can restrict
subsequent measurement or CSI acquisition to a smaller beam subset. If the
candidate set does not contain a sufficiently good beam, the selected beam and
the resulting CSI/PMI abstraction may be mismatched. We represent this effect
through a generic mismatch factor \(\eta_{b,u}(t)\), where larger values
correspond to more accurate beam and channel information.

For multi-user downlink evaluation, let \(\tilde{\mathbf{h}}_{b,u}(t)\) denote
the effective channel used by the evaluation module after the selected beam and
any mismatch abstraction have been applied. For a scheduled user set
\(\mathcal{U}^{\mathrm{sch}}_b(t) \subseteq \mathcal{U}_b(t)\), the BS forms a
ZF precoder
\begin{equation}
    \mathbf{W}_b(t)
    = [\mathbf{w}_{b,u}(t)]_{u \in \mathcal{U}^{\mathrm{sch}}_b(t)}
    \label{eq:zf-precoder}
\end{equation}
from the estimated or effective channel matrix. The resulting SINR of user \(u\)
is abstracted as
\begin{equation}
    \gamma^{\mathrm{ZF}}_{b,u}(t)
    =
    \frac{
    P_{b,u}(t)\left|\tilde{\mathbf{h}}_{b,u}^{\mathrm{H}}(t)
    \mathbf{w}_{b,u}(t)\right|^2
    }{
    \sum\limits_{v \in \mathcal{U}^{\mathrm{sch}}_b(t), v \neq u}
    P_{b,v}(t)\left|\tilde{\mathbf{h}}_{b,u}^{\mathrm{H}}(t)
    \mathbf{w}_{b,v}(t)\right|^2 + \sigma^2
    }.
    \label{eq:zf-sinr}
\end{equation}
Under perfect channel knowledge and feasible ZF, the inter-user interference
term is ideally suppressed. In the present abstraction, residual interference
may remain because sensing-assisted beam selection, target-user association, or
tracking errors can lead to imperfect effective-channel information.

The spectral efficiency of user \(u\) is then evaluated as
\begin{equation}
    r_{b,u}(t)=\log_2\left(1+\gamma^{\mathrm{ZF}}_{b,u}(t)\right).
    \label{eq:spectral-efficiency}
\end{equation}
Let \(\rho_{b,u}(t) \in [0,1]\) denote the fraction of the epoch consumed by
beam training, beam measurement, calibration, or re-association overhead for
user \(u\). The overhead-adjusted throughput is modeled as
\begin{equation}
    R^{\mathrm{adj}}_{b,u}(t)
    =
    \left(1-\rho_{b,u}(t)\right) B_b r_{b,u}(t),
    \label{eq:overhead-adjusted-throughput}
\end{equation}
where \(B_b\) is the bandwidth assigned by BS \(b\) in the simplified evaluation.
The system sum throughput at epoch \(t\) can be written as
\begin{equation}
    R^{\mathrm{sum}}_b(t)
    =
    \sum_{u \in \mathcal{U}^{\mathrm{sch}}_b(t)}
    R^{\mathrm{adj}}_{b,u}(t).
    \label{eq:sum-throughput}
\end{equation}
This abstraction allows the evaluation to account for two competing effects:
sensing-assisted prediction can reduce beam training overhead, while prediction,
association, or tracking failures can reduce the SINR through beam mismatch or
imperfect effective-channel estimation.

TODO[unstable-method]: Fix the mathematical mapping from Top-\(K\) beam miss,
beam mismatch, association failure, and tracking error to \(\eta_{b,u}(t)\) and
\(\tilde{\mathbf{h}}_{b,u}(t)\). Also fix the time or resource unit used to
compute \(\rho_{b,u}(t)\).

Similar system-level abstractions are commonly used in predictive beamforming
and ISAC studies to connect sensing, beam selection, and communication
performance through SINR, rate, or throughput-oriented metrics
\cite{liu2020radarassistedpredictivebeamformingvehicular,yuan2020bayesianpredictivebeamformingvehicular,liu2022learningbasedpredictivebeamformingintegrated,Wang_2023}.

\section{Problem Statement}
\label{sec:system-problem-statement}

Given the multimodal sensing observations
\(\{\mathcal{O}^{\mathrm{sen}}_b(t)\}_{t=0}^{T-1}\) and the communication
measurements available during initial access or later calibration, the system
aims to produce four coupled outputs. First, it detects vehicle targets in the
BEV frame. Second, it predicts a per-target beam power distribution and a
Top-\(K\) candidate beam set for each detected target. Third, it associates
detected sensing targets with communication user identities and maintains those
identities over time through tracking. Fourth, it evaluates the resulting
beam-management decisions through SINR, spectral efficiency, and
overhead-adjusted throughput under the simplified ZF-based communication model.

The problem can be summarized as the estimation and decision process
\begin{equation}
    \left\{\mathcal{O}^{\mathrm{sen}}_b(t),\mathbf{y}^{\mathrm{SSB}}_{b,u}(t_0),
    \mathcal{M}_{b,u}(t)\right\}
    \longrightarrow
    \left\{\mathcal{D}_b(t),\hat{\mathbf{g}}_{b,i}(t),
    \hat{\mathcal{K}}_{b,i}(t),a_{b,i,u}(t),
    \mathbf{x}_{b,r}(t),m_{b,u}(t)\right\},
    \label{eq:problem-mapping}
\end{equation}
where \(\mathcal{M}_{b,u}(t)\) denotes any later partial communication
measurement used for calibration or recovery. The desired behavior is to improve
beam prediction accuracy, maintain reliable target-user association and tracking,
reduce repeated full beam sweeps, and increase overhead-adjusted throughput. In
the current stage of the work, this objective is treated as a coupled
system-design and evaluation problem rather than as a single closed-form
optimization problem.

More formally, the final framework should balance the following goals:
accurate BEV target detection, high Top-\(K\) beam coverage, stable target-user
identity maintenance, low beam-management overhead, and high
\(R^{\mathrm{adj}}_{b,u}(t)\) under the evaluation model in
Section~\ref{sec:system-performance-abstraction}. The exact training losses,
matching costs, overhead accounting, and mismatch penalties will be specified
after the method and simulation pipeline are finalized.

TODO[unstable-method]: Convert this high-level problem statement into the final
mathematical formulation after the learning objectives, association cost,
tracking update, beam mismatch model, and evaluation protocol are fixed.
